\chapter{第十九讲\quad 财产犯罪}
\label{lecture:19}
\noindent{\small\hyperref[ref:lecture:19]{查阅本讲扩展规则}}\par

先固定行为人 $a$、原占有人 $b$、财物或利益 $x$ 和取财时点 $\tau$。写 $P_t(p,x)$ 表示 $p$ 在 $t$ 时实际或者按社会观念占有 $x$，$O(p,x)$ 表示所有权。$t^-,t^+$ 分别是某行为前后，$\Rightarrow$ 表示按所标规则推出，$\leadsto$ 表示案件中另有证据支持的因果联系。处在后一个时点的事实不能自动补足前一时点的故意。以下罪名判断均假定责任等共通条件已满足，明示数额条件的地方以经核验的适用标准为准。

\section{第一节\quad 基本原理}

\paragraph{教材模型}
令 $E(a,x,\tau)$ 为行为时有排除他人占有的意思，$U(a,x,\tau)$ 为使自己或第三人利用财物或其价值的意思。用 $I:=E\land U$ 表示非法占有目的的教学分解。这个模型解释取得与毁弃的分界。为保留款项才附条件返还财物，不能直接当成原物无条件归还。占有的存续也与所有权分开判断。教材比较「合法占有、所有权、平稳的非法占有与普通第三人」，这是不同主体之间的保护关系，不能把其中任何一项当作所有权凭证。

\paragraph{推导与反例}
甲将相机合法借给乙，期内偷偷从乙的住处取回并打算永久留用。题设给出 $O(a,x)$，也给出行为前 $P_{\tau^-}(b,x)$ 的合法依据。行为后甲排除乙使用并留为己用，故有 $E,U$，满足 $I$，可进入盗窃审查。只凭 $O(a,x)$ 推不出 $\neg I$。若改为乙刚抢走相机，甲在侵害尚进行时依法当场夺回，合法占有和权利恢复的前提均变了，须另查防卫等。若甲拿走财物只为毁弃，通常缺 $U$，转查毁坏型罪名。电力、虚拟财产、债权凭证等须先证明其法律评价中的财产价值和控制方式，违禁物的刑法保护及计价见第十一节。

\section{第二节\quad 侵占罪}

\paragraph{法条与分解}
第270条有两条入口。$T$ 为行为人已独立占有代为保管的他人财物，$L$ 为他人的遗忘物或者埋藏物已经进入行为人占有，$M$ 为非法据为己有，$Q$ 为数额较大，$R_T,R_L$ 分别为拒不退还和拒不交出。检验式是
\[
 M\land Q\land[(T\land R_T)\lor(L\land R_L)] .
\]
告诉才处理是该条的程序条件。教材对出售、消费等处分能否表达拒不退还另有讨论，稳妥做法是把 $M$ 和相应的拒不交还事实分别列明。对不法原因委托及返还请求，须标明所采解释。

\paragraph{推导与反例}
甲收取乙委托保管的贵重相机，在自己家独立保管。后来甲擅自卖掉，乙索还时甲明确拒绝，数额已达标准。先前取得占有在合法委托时完成，故有 $T$。出售给出 $M$，拒绝给出 $R_T$，连同 $Q$ 得第270条路径。若乙只在桌边借给甲查看，仍保持现实控制，甲偷带离开，则此前不能确立 $T$，应查盗窃所需的占有转移。路旁遗失物若仍处于场所管理者或原主的规范占有中，亦不可仅凭「捡到」推出 $L$。

\section{第三节\quad 盗窃罪}

\paragraph{法条与教材模型}
记 $T(a,b,x,\tau)$ 为行为人违反原占有人意思平和地打破其占有并建立自己或第三人占有的行为。教材采用个别公开说时，平和性与秘密性分开。刑法第264条规定数额与四个特殊入口的析取。设 $Q$ 为数额较大，$M,H,W,P$ 分别为多次、入户、携带凶器、扒窃，则
\[
 I\land T\land(Q\lor M\lor H\lor W\lor P)
 \quad\Longrightarrow\quad \text{盗窃既遂审查}.
\]
这里 $T$ 取完成的转移。仅着手而未转移的，另查第23条和盗窃司法解释规定的未遂条件。两高盗窃解释第3条分别界定二年内三次、非法进入供家庭生活且与外界相对隔离的住所、携带有关器械、公共场所或公共交通工具内窃取随身携带财物。

\paragraph{推导与反例}
甲在公交车上徒手从乙随身背包取出钱包，乙未同意，甲取得控制并打算占有。原占有由乙对背包的随身控制得出，取出与控制建立 $T$，目的给出 $I$，地点与对象推出 $P$。由 $P$ 推出析取式，故不必再以 $Q$ 为唯一入罪入口。若钱包在乙身旁的桌面而非随身携带，不能沿用 $P$。若主人目睹甲平和搬走财物但无法阻止，个别公开说仍允许 $T$，全部秘密说则另加秘密性条件，答案须注明理论前提。以为只要公开便成立抢夺，会跳过抢夺所需的对物夺取。

\section{第四节\quad 抢夺罪}

\paragraph{法条与教材模型}
教材用 $F_x$ 表示乘隙对被紧密控制的财物施力夺取，$F_p$ 表示为压制反抗而对人实施强制。依第267条第一款，抢夺入口是 $I\land F_x\land(Q\lor M)$，$Q,M$ 为法定数额或多次条件。用 $W$ 表示依法认定的携带凶器。$W\land F_x$ 进入第二款按抢劫定罪的路径。若以器械威胁压制反抗，优先检验第263条直接抢劫。依抢夺司法解释第6条，驾驶车辆取财的三类具体情形还要单列审查。

\paragraph{推导与反例}
甲突然扯走乙手中紧握的手机，手段直接作用于手机，无对人暴力或威胁，数额达标且未携带凶器。事实给出 $F_x,I,Q,\neg F_p,\neg W$，由第一款进入抢夺。若甲先以暴力逼乙放手再拿手机，则因果图改为 $F_p\leadsto\text{不能反抗}\leadsto\text{取财}$，应查抢劫。若甲仅平和拿走旁边桌上无人紧握的财物，$F_x$ 缺失，不能因为当场被看见就称抢夺。

\section{第五节\quad 抢劫罪}

\paragraph{第263条的直接路径}
记 $F$ 为为取财而实施的暴力、胁迫或其他压制手段，$S$ 为对占有人或可阻止取财的人压制反抗，$G$ 为取得财物。教学上的完成图为
\[
 F\leadsto S\leadsto G .
\]
图上每条边均需案件事实，先前的暴力与后来的取财偶然相邻不能代替两条边。按照最高法《两抢意见》确立并由抢劫审判文件衔接的标准，在抢劫构成前提下，取得财物或造成被害人轻伤以上后果之一，即可按既遂评价。故判断既遂不能只写 $G$。第263条的直接抢劫通常不以数额较大为入罪条件。升格应逐项按八项法定情形检查，如「入户抢劫」查户的性质、入户目的及抢劫发生地点，「持枪抢劫」与第267条的携带凶器拟制路径分别判断。

\paragraph{第269条的转化路径}
设 $A$ 为作为前提的盗窃、诈骗、抢夺，$P$ 为窝藏赃物、抗拒抓捕、毁灭罪证任一目的，$L$ 为当场，$F'$ 为使用暴力或以暴力相威胁。转化规则为
\[
 A\land P\land L\land F'\Longrightarrow
 \text{依第263条定罪处罚}.
\]
依2016年抢劫指导意见，$A$ 一般只须前行为已经着手。财物数额明显低于数额较大标准且无《两抢意见》第五条列举的五种情节时，不构成转化型抢劫。以强度较小的摆脱方式逃离，且未造成轻伤以上后果的，可以不认定为使用暴力。「当场」包括作案现场以及刚离现场即被发现并追捕等连续情形。对共同实施前行为者，另证对转化暴力的共同故意和参与。

\paragraph{推导与反例}
甲为取得乙的包，先以暴力压制乙反抗，因而取得包。目的、压制、取财及两条因果边均由题设给出，故第263条直接抢劫既遂。如果暴力前已经以盗窃取得包，甲在连续追捕中为了抗拒抓捕实施足以认定的暴力，则代入 $A,P,L,F'$ 后走第269条。若第二天另遇失主，甲出于报复殴打乙，既缺 $L$ 又缺三类目的 $P$，保留前次盗窃，后次暴力另论。若甲用人质向另一人提出不法要求，则另查绑架，不能仅因出现强制和钱就套用抢劫。

\section{第六节\quad 敲诈勒索罪}

\paragraph{法条与模型}
第274条以数额较大或者多次敲诈勒索等为入口。教学模型在 $I$ 下写作
\[
 \text{威胁或要挟}\leadsto\text{恐惧}
 \leadsto\text{基于恐惧处分}\leadsto\text{取得与损失}.
\]
抢劫与敲诈都可能有威胁。检查强制的现实、紧迫程度、能否压制反抗以及财物如何取得。「当场给钱」可作为证据，单独不能推出抢劫。恐惧与认识错误可能并存，要证明哪一项实际导致处分。对于维权索财，实际债权、主张范围和非法占有目的各查一次，手段不法也可能构成其他犯罪。

\paragraph{推导与反例}
甲为非法获利威胁乙，称若三日内不给钱就公开乙的隐私。乙因恐惧转账，甲收到，数额达标。证据依次建立威胁、恐惧、处分、取得和损失，且不存在当场压制反抗的事实，故按第274条推导敲诈勒索既遂。若乙识破虚张声势而并未付款，处分与取得均缺失，改查着手及未遂。若甲主张的是确有依据的本人债权，仍需查手段责任，但不能从债权人言辞强硬直接推出 $I$。

\section{第七节\quad 诈骗罪}

\paragraph{第266条的链式模型}
用 $B$（欺骗）、$E$（认识错误）、$D$（合格处分）、$G$（行为人或其指定人取得）与 $L$（损失）建立以下结构：
\[
 B\leadsto E\leadsto D,\qquad
 D\leadsto G,\qquad D\leadsto L .
\]
另需非法占有目的、适用数额等入罪条件，且没有优先适用的特别诈骗条款。每个节点存在、每条边成立，是两组独立命题。对积极的虚构、隐瞒和不作为欺骗分别审查，尤其要为不作为指出告知义务。不能由合同落空或最终亏损反推最初欺骗。

\paragraph{两角案件的两次验算}
把 $D$ 分解为 $D_o$（直接改变占有或财产状态）和 $D_s$（意识到正处分相应财物）。这是教材采用的处分意识必要说。若甲调换标价，收银员看到具体相机并以错误价格主动交出，则认识错误是价格而非相机是否存在。交付同时给出 $D_o,D_s$，价差损失有证据时可按诈骗推导。若甲把相机藏在饼干盒，收银员只意识到处分饼干，依本说，相机缺 $D_s$。诈骗链在 $D$ 处断开，改查甲打破相机占有的盗窃。整体计量的数量错误与按件加塞不同。把同类物作为一整袋交付时，重量错误可能仍有整体处分意识。按件交付而对被加塞的额外一件毫无认识时，须对额外物单独判断。另有处分意识不要说，结论分歧来自 $D$ 的定义变化。

以调换标价为例，假定价差达到适用数额，证据集 $\Gamma$ 包含调换动作、收银员的误信、对相机的认识和交付、少收货款等事实。以 $\vdash$ 表示由这些事实和本节采用的规则推出，写出三步：
\[
\begin{aligned}
 \Gamma&\vdash B,\ E,\ B\leadsto E,\\
 \Gamma&\vdash D_o,\ D_s,\ E\leadsto D
           \quad\Longrightarrow\quad \Gamma\vdash D,\\
 \Gamma&\vdash D\leadsto G,\ D\leadsto L,\ G,\ L,\ I,\ Q
           \quad\Longrightarrow\quad \Gamma\vdash\text{诈骗既遂}.
\end{aligned}
\]
第三步使用前两步已经证明的命题。若独立证据证明收银员发现标价错误，仍依据获授权的促销价主动出售，$B$ 可存在，但 $B\leadsto E$ 及 $E\leadsto D$ 均失效。价格变化本身不是因果证据。

\paragraph{直接性反例}
甲骗称乙的孩子在楼下出事，乙下楼查看，甲随后独自进入乙家拿走相机。谎言导致乙离开，但离开只为取财创造条件。若删除甲后来拿相机的行为，占有并未因乙下楼而转移，故 $\neg D_o$，不能推出诈骗。恰因占有变化由后续取走实现，转查盗窃。这一步的删去检验可复用到魔术演示、借手机和误认保管的题目。

\paragraph{三角案件的权限验算}
设被骗人 $v$、受损人 $w$。$v$ 可与 $w$ 不同，但要证明被骗人 $v$ 的错误导致其处分，并且 $v$ 对 $w$ 的财产具有刑法评价所需的处分地位 $A(v,w,x)$。教材列出关于处分地位范围的不同见解。若甲以假提货单欺骗受公司授权核单交货的仓管员，仓管员因错误将相应商品交给甲，且公司遭受损失，则有 $A,D_o,D_s$，诈骗链闭合。若甲骗不相干的路人代拿别人锁在仓库的财物，则该人仅事实接触而无处分地位，转查利用他人实施盗窃的间接正犯。两种结论不能由「经过第三人交付」一语分流。

\section{第八节\quad 职务侵占罪}

\paragraph{法条与模型}
第271条第一款的检查式为
\[
 \text{单位工作人员}\land\text{职务便利}
 \land\text{本单位财物}\land\text{实际非法占为己有}
 \land I\land Q .
\]
职务便利须能说明行为人因职务而主管、管理、经手财物。对本单位财物的外延，另审查单位所有或依法管理的具体关系。公式用于既遂路径，未得逞的另审形态。依第二款，对国有单位中从事公务及受国有单位委派从事公务的相关人员，符合本条行为的，依法转向第382、383条。按法释〔2026〕6号第8条，第271条的定罪量刑标准参照贪污罪，仍须综合犯罪性质、情节与社会危害性。

\paragraph{推导与反例}
甲为非国有单位仓库管理员，利用本人保管权限将单位库存据为己有，数额达标。身份文件给出工作人员身份，授权和取货方式证明职务便利，库存归属证明对象，处分方式和目的证明 $I$，由合取式进入第271条。若甲是同公司的普通职工，只知道库存的位置，趁人不备破锁偷取，则「知道地点」仅是工作便利的线索，不能证明职务上管理或经手权限，转查盗窃。若只有临时使用意思，则另查第272条，不能用事后未还直接替代 $I$。

\section{第九节\quad 普通罪名}

\paragraph{法条与分类}
第275条查故意毁坏财物及数额较大或其他严重情节。教材把物理破坏、价值或功能破坏分开。第272条查职务便利下挪用本单位资金归个人使用或借贷他人，并满足法定的用途、数额和时间组合。其三条分支分别是数额较大且超过三个月未还、数额较大且营利、进行非法活动，具体司法适用条件依有效规范核验。第273条则查特定款物用途被擅改、情节严重、国家和人民群众利益重大损害。第276条之一可拆作
\[
 (\text{转财或逃匿以逃付}\lor\text{有能力而不付})
 \land Q\land\text{政府有关部门责令后仍不付}.
\]

\paragraph{推导与反例}
甲经营企业，拖欠劳动报酬达到适用数额且有能力支付，政府有关部门依法责令后仍拒付。先由支付能力与不支付得第二支，再合取 $Q$ 和责令后仍不支付，推出第276条之一的基本路径。若只证普通拖欠，且尚无政府有关部门责令，最后一项不成立，不能推出该罪。另如单位人员擅改救灾款专用用途并造成法定重大损害，首先代入第273条，而非把「临时挪走」直接等同第272条的归个人使用。

\section{第十节\quad 财产犯罪的既遂}

\paragraph{教材模型}
对常见有形物取得犯，设 $L_b$ 为原占有人失去有效控制，$G_a$ 为行为人或其指定人建立有效控制，取财结果用 $K:=L_b\land G_a$ 检验。地理移动、过门槛、短时间接触均不是一般充分条件。诈骗和敲诈还须验证处分引发损失和取得。抢劫的轻伤以上替代路径见第五节。教材另有不同控制观，写题时应指出题设和采纳的规则。

\paragraph{推导与反例}
甲将乙的珠宝放进本人专属、已上锁的储物格，乙对格内财物已经不能控制，甲独有开锁条件。题设据此证明 $L_b,G_a$，故得 $K$。之后立即被抓不抹去已完成的控制转移。反之，甲才从柜台拿起珠宝即被店员夺回，尚未排除柜台方控制，$L_b$ 未立，不能以曾经接触推出既遂。乙因被骗把款汇给与甲无关且甲不能支配的第三人，则乙虽可能有损失，甲的取得缺失，仍须单查诈骗未遂。敲诈款仅放在公共场所约定位置时，教材分别用行为人支配程度及原主可取回程度论证，不能以到达坐标代替控制事实。

\section{第十一节\quad 财产犯罪的数额}

\paragraph{法条、司法解释与模型}
把罪名 $c$、地域 $j$、时点 $t$ 及具体情节 $h$ 写进门槛函数 $q(c,j,t,h)$。客观价值 $v$、主观认识的价值范围 $J$ 和未得部分分别记账。教学上的主客观相一致检查是：适用数额档时，先证客观达到，再审查行为人对对象与价值情状的认识能否覆盖该档，不能以一个精确报价代替主观故意。对盗窃等特殊入口及未遂须另按规定，不能一律要求完成数额达到门槛。盗窃违禁品依两高盗窃解释第1条，按照盗窃罪处理的，根据情节轻重量刑，不能直接采用票面数字或把估算的黑市价格当统一的法定计价标准。

\paragraph{推导与反例}
某合同诈骗案题设已认定甲意图骗取 $100$ 万元，已实际控制被害人支付的 $30$ 万元，尚有 $70$ 万元部分符合已着手未取得的条件。设计划额 $V_p$、取得额 $V_c$，则在同一对象范围、无重复计入时
\[
 V_u=V_p-V_c=100-30=70\text{ 万元}.
\]
应分别检查两部分对应的定性、量刑幅度和适用规则。不能把未遂写成 $100$ 万元再另加既遂 $30$ 万元。若剩余款项尚未进入实行阶段，则等式只给出「计划而未得」的差额，不能单凭差额推出未遂。依盗窃解释第12条，盗窃既有既遂又有未遂，分别达到不同量刑幅度的取处罚较重的规定，达到同一幅度时按盗窃既遂处罚，其他罪名须核验各自规则。

\section{第十二节\quad 财产性利益}

\paragraph{教材模型与争议}
设 $s_w$ 为被害人对某项具体利益在案发前后的可实现状态，$s_a$ 为行为人取得的经济利益。若欺骗或强制使可依法主张的债权实际被处分，按教材所列多数说可以评价财产性利益。少数说对「转移」要求更严格。违法债权是否受保护也有分歧。适用多数说，须证明法秩序认可和经济价值两项，不能把民事债务存在、迟延支付和实际免除混写。
\[
 \text{已有合法可实现债权}\land
 \bigl(\text{欺骗}\leadsto\text{权利人处分}\bigr)
 \land\text{损失与对应收益}
 \Longrightarrow\text{审查利益诈骗}.
\]

\paragraph{推导与反例}
甲本来打算支付服务费。服务完成后，甲编造已结清的凭证，具有结算权限的债权人因误信作出具有实质处分效果的决定。若题设证明债权的可实现性实际受损、甲取得对应利益且数额达标，依多数说，诈骗链发生在财产性利益上。若甲只悄悄离开，债权仍可正常主张，则欠费存在不能推出处分和损失，教材对该类逃费是否侵害可处罚的财产性利益另列争议。若从起初即欺骗店家交付饭菜，须先评价饭菜的占有转移，不要只盯住后来形成的价款债权。

\paragraph{账户与服务}
将账户标识、取款或支付权限、存款债权和资金实际归属画成四列。知道口令不等于拥有款项；错汇、卖卡后划转、欺骗柜员和通过机器取款，也应按授权及实际转移方式分别检验。例：甲收到误汇通知后，题设仅给出收款事实，不能推出其在汇款前作了欺骗。之后拒还的民事返还义务和是否满足具体犯罪条件另行评价。第265条对以牟利为目的的盗接通信线路、复制电信码号等有专门指向。其他电、电话、网络利益还需核实具体取得方式与计量依据。

\section{第十三节\quad 财产犯罪的罪数联系}

\paragraph{教材模型}
用事件序列 $(a_1,a_2)$ 记前后行为，对每一项写「对象、受害人、手段、损失、法条」。记 $\Delta_i$ 为第 $i$ 步可归责的财产损害。如果 $a_2$ 只是把前罪所得变现，不能只因又发生交易就加一个取得罪。如果 $a_2$ 另有相对人、欺骗与处分，并造成独立损失，则重新计算第二条犯罪链。最后检验法定的并罚、转化规则和不可罚的事后行为。教材对同一份财产损害能否分别评价列有不同处理，不能以「动作数等于罪数」代替论证。

\paragraph{推导与反例}
甲盗得乙的相机，之后虚构自己有权处分，骗丙支付购机款。题设明确丙未取得相机或等价财产。先由乙的占有转移得盗窃，损害为 $\Delta_1$。再由欺骗丙、丙的误信及交款、甲取得和丙损失得诈骗，损害为独立的 $\Delta_2$。在其余条件成立且无特别规则时分别评价。若丙明知是赃物仍协助变现，第二条链缺认识错误，应查丙的掩饰、隐瞒犯罪所得等责任，而非照抄诈骗。若是原受害人因再次欺骗免除返还义务，要先证新增、可分的财产性利益损害，警惕对同一损失重复归责。先偷后以虚假退货换取价款的案件也按这一原则比较。

\section*{闭卷记忆与验算}
\paragraph{复述路线}
写「占、转、手、意、果、额、并」。占是 $P_{\tau^-}$ 与 $O$ 分离。转是问财物由谁移转，行为人自己打破占有，还是被害人基于错误或恐惧处分。手是平和、对物夺取、对人压制。意是 $I=E\land U$ 与各罪特别目的。果是实际控制、损失及归责。额是入罪门槛、特殊入口、既遂与未遂分账。并是先判每个行为，再检验转化、特别条款和重复评价。复习诈骗时另写五字「骗、错、处、得、损」，再给「处」补上直接性、处分意识和处分地位三问。

\begin{enumerate}
\item 题：店员交付纸箱，未察觉其内被藏了一块手表。按处分意识必要说，何处断链？答：对手表有物理交付，但缺对其存在与处分的认识，$D_s$ 缺失，按教材模型转查盗窃。若店员看见手表，只被虚假价格欺骗，则 $D_s$ 可成立。
\item 题：行为人取得某物后，在连续追捕中为毁灭证据以暴力相威胁。还须核验什么？答：前行为是否属于第269条的盗窃、诈骗、抢夺，现场延续性、暴力威胁的法定程度及故意。证齐才进入转化路径，之后另判既遂、未遂和升格。
\item 题：将钱包拿起来马上被追回，能否直接判既遂？答：先看原主是否失去有效控制、行为人是否建立控制，均不能由一次触碰推出。若已经装入行为人专属空间并排除原主支配，立刻追回不当然推翻既遂。
\end{enumerate}
